\documentclass[12pt]{article}
\usepackage[utf8]{inputenc}
\usepackage[T1]{fontenc}
\usepackage{hyperref}
\title{Mbandja Language Resources for K-12 Curriculum Development}
\author{Compiled by Claude Sonnet 5.5 (claude-sonnet-5-5) for Andr\'as Kornai}
\date{}
\begin{document}
\maketitle

\section{Introduction}
Mbandja (also Banja, Mbanza; ISO 639-3 zmz, Glottocode mban1263) is the largest of the Banda languages, spoken mainly in the Democratic Republic of the Congo, with smaller communities in the Republic of Congo and the Central African Republic. Published speaker estimates vary widely between sources. Web searches found very little material written for schools. This report lists only what appeared in search results.

\section{Primary Educational Resources}
No primer, textbook or school curriculum for Mbandja turned up in search results. Nothing found indicates that the language is taught in schools.

\section{Linguistic and Reference Materials}
\begin{itemize}
\item Wikipedia article on the Mbandja language, with phonology and dialect notes (Kala, Gbado) and a citation of a 1978 grammatical study of the Balawo dialect by Nyi-Zonga Tingbo: \url{https://en.wikipedia.org/wiki/Mbandja_language}
\item Glottolog entry for Mbandja: \url{https://glottolog.org/resource/languoid/id/mban1263.bigmap.html}
\item Numeral system dataset for Mbanza (Mabandja), DR Congo, from the Max Planck Institute for Evolutionary Anthropology: \url{https://lingweb.eva.mpg.de/channumerals/Mbanza.htm}
\end{itemize}

\section{Religious and Community Resources}
\begin{itemize}
\item Joshua Project profile of the Mbanza people in DR Congo: \url{https://www.joshuaproject.net/index.php/people_groups/13620/CG}
\item Joshua Project profile of the Mbanza people in the Republic of Congo: \url{https://joshuaproject.net/people_groups/13620/CF}
\end{itemize}
No Mbandja Bible translation was found in the search results.

\section{Digital and Audio Resources}
Global Recordings Network has a page for code zmz: \url{https://glorecnet-website.k.globalrecordings.net/en/language/zmz}. According to the search summary, recordings exist in the related Gbado and Kala varieties, but none were listed for Mbandja itself.

\section{Curriculum Development Context}
Very little documentation was found, and no school materials. A curriculum would have to start with a standardized orthography, a lexicon, and community consultation. Lingala and French are the likely languages of wider communication, and neighbouring Banda materials may offer a model. All the items above are reference or background material, not classroom resources.

\end{document}